\documentclass[12pt]{article}
\usepackage[utf8]{inputenc}
\usepackage[round,authoryear]{natbib}
\usepackage{graphicx}
\usepackage{amsmath}
\usepackage{amssymb}
\usepackage{times}
\usepackage{booktabs}
\usepackage{algorithm}
\usepackage{algorithmic}
\usepackage{geometry}
\geometry{
 a4paper,
 total={170mm,257mm},
 left=20mm,
 top=20mm,
}

% headline number filled in from experimental_results/NUMBERS.json
\newcommand{\forecastmse}{0.044}

\title{Causal discovery Kolmogorov--Arnold networks for interpretable time-series forecasting in managerial decision support}
\author{Quang-Vinh Dang \and Dat Le \and Minh Ngoc Dinh}
\date{\today}

\begin{document}

\maketitle

\begin{abstract}
Understanding causal relationships in organizational time-series data is critical
for managerial decision-making, yet forecasting systems often trade
interpretability for accuracy. We introduce \textbf{Causal Discovery
Kolmogorov--Arnold Networks (CD-KAN)}, an information-systems artifact that, from
a single differentiable model, jointly produces a one-step forecast and an
interpretable, lag-resolved causal graph. CD-KAN predicts each variable \emph{only} through per-(cause, lag) B-spline edge
functions of its candidate parents---an information bottleneck that makes the
learned structure identifiable---with an edge-level group-lasso for sparsity and
an optional acyclicity constraint on contemporaneous edges. We evaluate it
transparently against seven genuine baselines (PCMCI, VAR-LiNGAM, NOTEARS, GOLEM,
VAR-Lasso, and two KAN-based Granger variants) using threshold-free AUROC/AUPRC
with paired significance tests. On synthetic data with known ground truth, CD-KAN
recovers linear structure perfectly (AUROC~$=$~1.0) and is competitive on
non-linear structure (AUROC~$=$~0.82 overall, $\approx$0.95 at $d{=}5$),
significantly outperforming the KAN-Granger baselines ($p<0.05$) and matching
several classical methods, while remaining fully interpretable. We show that this
identifiability hinges on the bottleneck: a variant with a dense forecasting
backbone collapses to chance. Under a leakage-free rolling-origin protocol CD-KAN
also forecasts an 8-asset financial panel competitively (MSE $\approx\forecastmse$;
we retract an earlier $0.0008$ traced to a test-leaking normalization), its
learned real-market graph is temporally stable roughly six-fold above a
permutation chance level ($p<0.0001$), and its spline edge functions expose
linear, saturating, and threshold-like dependencies---supporting a ``glass-box''
decision-support tool for moderate-dimensional, non-linear business systems.
\end{abstract}

% Include the Introduction Section
% Include the Introduction Section
\input{jmis_introduction}

% Include the Method Section
% Include the Method Section
\input{jmis_method}

% Include the Results Section
% Include the Results Section
\input{jmis_results}

% Include the Discussion Section
% Include the Discussion Section
\input{jmis_discussion}

% Include the Conclusion Section
\input{jmis_conclusion}

\bibliographystyle{apalike}
\bibliography{references}

\end{document}
